\section{Analysis of Prior and Posterior Signals}

\noindent\textbf{Prior signals.} Prior signals are derived from agents’ initial responses. They typically measure output quality and are used to exclude agents whose initial responses already appear reliable from debate, and/or prune the communication graph to prioritize deliberation with high-quality agents, aiming to improve accuracy with less communication.

We study three prior signals: min. log-likelihood (min-LL), which captures the least likely token; perplexity (PPL), which measures response-level generation uncertainty; and the LLM's self-reported confidence (conf).

\noindent\textbf{Our proposal: SVR.} SVR measures how often an agent retains its initial belief after debating with peers. This evidence-grounded signal provides a posterior estimate of agent reliability that is less dependent on potentially hallucinated prior information. Let $D$, $r$, and $c$ denote the total number of debates involving an agent, the number of times the agent retains or changes its belief after debate, respectively. We define $SVR=(r-c) / D$.

\begin{figure}
\centering
\includegraphics[width=\columnwidth]{figures/motivation}
\caption{\label{fig:motivation}\textbf{SVR better identifies correct agents than prior signals under increasing difficulty.} Percentage of problems where the top-ranked agent under each signal is correct, as a function of increasing difficulty.}
\end{figure}

\noindent\textbf{Approach.} We validate SVR on two LLMs: GPT-OSS-120B~\citep{Openai2025gptoss} and DeepSeek-V3.1~\citep{Deepseek2024Deepseek}, using questions from IMO-AnswerBench~\citep{Luong2025Towards}. For each question, we run six agents and group questions by difficulty, measured by the number of agents that are correct before debate. We compare three prior signals against SVR by measuring how often the top-ranked agent under each signal is correct.

\noindent\textbf{Results.} Figure~\ref{fig:motivation} reports this percentage across difficulty groups, showing how each signal aligns with agent correctness as problems become harder. From 4+ correct agents (easiest) to 1 correct agent (hardest), the correctness of top-ranked agents under min-LL, PPL, and confidence drops from 67--87\% to 5--34\%. \textbf{This sharp decline shows that prior signals become unreliable on difficult problems, where hallucinated or overconfident reasoning makes model behavior poorly aligned with correctness.} In contrast, SVR remains predictive as difficulty increases: among questions with only 2 and 1 correct agents, SVR achieves 68--78\% and 47--49\%, respectively, outperforming the strongest prior signal by 18--23 and 15--22 points.

% While SVR is reliable, it is only observable after debate. Next, we present \sysname, an efficient MAD algorithm that incrementally estimates SVR from early communications and uses it to guide subsequent communication.
